% !TEX root = ../../../main.tex
% 所属：第一部分：单变量分析学 / 第5章：级数理论
% 节：5.1 级数的定义
% 正文按原文件顺序保存；由 ../chapter5.tex 调用。
\section{级数的定义}
\begin{definition}
	形如 $a_1 + \dots + a_n + \dots$（$a_1, \dots, a_n, \dots \in \mathbb{R}$）叫作数项级数，记为 $\sum\limits_{n=1}^{\infty} a_n$，简写为 $\sum a_n$。
	
	由 $\sum a_n$ 可以构造一个部分和数列 $\{S_n\}_{n=1}^{\infty}$：
	\[
	S_n = \sum_{i=1}^{n} a_i
	\]
	
	称 $\sum a_n$ 收敛，当且仅当 $\{S_n\}$ 有一个有限的极限 $S$，称 $\sum a_n$ 收敛于 $S$，记为
	\[
	\sum_{n=1}^{\infty} a_n = S
	\]
	
\end{definition}

注意：不能通过 $S_n$ 的通项直接研究 $\sum a_n$，一般而言 $S_n$ 的通项难以求得，因此我们要从 $a_n$ 的角度研究 $\sum a_n$ 的收敛性。


\begin{example}
	\begin{enumerate}
		\item 等比级数 $1 + q + q^2 + \dots + q^n + \dots$
		
		部分和：
		\[
		S_n = 1 + \dots + q^{n-1} = \frac{1 - q^n}{1 - q}
		\]
		\begin{itemize}
			\item 当 $q = 1$ 时，$S_n = n \to +\infty$，级数发散；
			\item 当 $q = -1$ 时，$S_n$ 在 $0$ 和 $1$ 之间振荡，级数发散；
			\item 当 $|q| < 1$ 时，$S_n \to \dfrac{1}{1-q}$，级数收敛；
			\item 当 $|q| > 1$ 时，$S_n \to \infty$，级数发散。
		\end{itemize}
		\item $p$-级数
		\[
		\sum_{n=1}^{\infty} \frac{1}{n^p} = 1 + \frac{1}{2^p} + \dots + \frac{1}{n^p} + \dots
		\]
		
		部分和：
		\[
		S_n = 1 + \frac{1}{2^p} + \dots + \frac{1}{n^p}
		\]
		
		由 \ref{调和级数} 可知：
		\begin{itemize}
			\item 当 $p > 1$ 时，$\sum \dfrac{1}{n^p}$ 收敛；
			\item 当 $p \le 1$ 时，$\sum \dfrac{1}{n^p}$ 发散。
		\end{itemize}

	\end{enumerate}

\end{example}

\begin{proposition}(收敛的必要条件)
	若级数 $\sum a_n$ 收敛，则必有
	\[
	\lim_{n \to \infty} a_n = 0
	\]
	
\end{proposition}
\begin{proof}
	\[
	a_n = S_n - S_{n-1}
	\]
	由于 $\sum a_n$ 收敛，故部分和数列 $\{S_n\}$ 收敛，设
	\[
	\lim_{n \to \infty} S_n = S
	\]
	则同时有 $\lim\limits_{n \to \infty} S_{n-1} = S$，因此
	\[
	\lim_{n \to \infty} a_n = \lim_{n \to \infty} (S_n - S_{n-1}) = S - S = 0
	\]
\end{proof}
\begin{proposition}
	若级数 $\sum a_n$ 与 $\sum b_n$ 收敛，则对任意 $x,y \in \mathbb{R}$，级数 $\sum (x a_n + y b_n)$ 收敛，且
	\[
	\sum (x a_n + y b_n) = x \sum a_n + y \sum b_n
	\]
	
\end{proposition}

当 $\sum a_n$ 与 $\sum b_n$ 至少有一个发散时，上述结论不一定成立,例如：
 \begin{enumerate}
 	\item $a_n = \dfrac{1}{n}$，$b_n = -\dfrac{1}{n} + \dfrac{1}{n^2}$，此时 $\sum a_n$ 发散，$\sum b_n$ 发散，但 $\sum (a_n + b_n) = \sum \dfrac{1}{n^2}$ 收敛。
 	\item 又如：$a_n = \dfrac{1}{n}$，$b_n = \dfrac{1}{n}$，此时 $\sum a_n$、$\sum b_n$ 均发散，且 $\sum (a_n + b_n) = \sum \dfrac{2}{n}$ 也发散。
 	\item 若 $\sum a_n$ 发散，$\sum b_n$ 收敛，则 $\sum c_n = \sum (a_n + b_n)$ 必发散。（证明：用反证法。若 $\sum c_n$ 收敛，则 $\sum a_n = \sum c_n - \sum b_n$ 也收敛，与题设矛盾。）
 \end{enumerate}
\begin{proof}
	设 $S_n = \sum_{i=1}^{n} a_i$，$T_n = \sum_{i=1}^{n} b_i$，由 $\sum a_n$、$\sum b_n$ 收敛，可知
	\[
	\lim_{n \to \infty} S_n = S, \quad \lim_{n \to \infty} T_n = T
	\]
	而级数 $\sum (x a_n + y b_n)$ 的部分和为 $x S_n + y T_n$，于是
	\[
	\lim_{n \to \infty} (x S_n + y T_n) = xS + yT
	\]
	故 $\sum (x a_n + y b_n)$ 收敛，且结论得证。
\end{proof}


我们考虑级数上对加法的结合律与交换律。

对结合律：仅对收敛级数成立，发散级数不成立。

对交换律：对收敛级数也不一定成立。

例：对发散级数 $\sum (-1)^n$
\begin{align*}
	(-1+1)+(-1+1)+\dots &= 0+0+\dots = 0 \quad \text{收敛} \\
	-1+(1-1)+(1-1)+\dots &= -1+0+\dots = -1 \quad \text{收敛} \\
	(-1+1-1)+(1-1+1)+\dots &= -1+1-1+\dots \quad \text{发散}
\end{align*}
关于交换律的问题，我们后续再说。
\begin{proposition}
	设级数 $\sum a_n$ 收敛，则在其表达式 $a_1 + a_2 + \dots + a_n + \dots$ 中，不改变求和次序，任意添加括号（每个括号内为有限项），所得新级数也收敛，且其和与原级数相同。
\end{proposition}
\begin{proof}
	记原级数的部分和为 $S_n = a_1 + a_2 + \dots + a_n$。
	
	将原级数添加括号后，可表示为：
	\[
	\sum a_n = (a_1 + \dots + a_{n_1}) + (a_{n_1+1} + \dots + a_{n_2}) + \dots = b_1 + b_2 + \dots + b_k + \dots
	\]
	其中 $b_k = a_{n_{k-1}+1} + \dots + a_{n_k}$，约定 $n_0 = 0$。
	
	设新级数的部分和为 $T_k = b_1 + b_2 + \dots + b_k$，则有
	\[
	T_k = b_1 + \dots + b_k = S_{n_k}
	\]
	即 $\{T_k\}$ 是原级数部分和数列 $\{S_n\}$ 的一个子列。由于原级数收敛，$\{S_n\}$ 收敛，因此其子列 $\{T_k\}$ 也收敛，且极限相同。
	
	这就证明了新级数收敛，且和与原级数相同。
	
\end{proof}

